\section{Cross-region rotation}
\label{sec:xprotocols}

\subsection{Four-module cross-region architecture}
\label{sec:runtime-modules}
Each correct regional replica runs a \defterm{Session Manager}. The
manager observes finalized L1 and Espresso history, receives regional
results, and manages independent programs named by $(\beta,r)$.
It is local replicated logic, not an additional trusted regional
controller. A region can run an old session while preparing a possible
successor and handing an earlier session's data duties to storage and
publication workers. The full regional view remains $(\beta,r,g,t)$.

The cross-region interface has four modules. The Soft Force Detector
recomputes delivery failure from authenticated history. The Multisig
Health Controller submits the existing L1 controls from health and
operational observations. These are parallel sources of intervention.
The Cross-region Handler manages candidate and fixed handoffs. Each
session's Ready and Activation Gate checks its exact authorized starting
point before entering the regional engine.

\begin{table}[htbp]\centering\small
\caption{The four cross-region modules.}
\label{tab:runtime-modules}
\begin{tabularx}{\linewidth}{@{}p{.22\linewidth}Y Y@{}}\toprule
Module and executor & Inputs & Output and authority limit\\\midrule
Soft detector; any verifier & Complete verified Espresso history,
export duty, fixed windows and accepted counts & Recomputable failed-window
condition; it does not itself change $\beta$.\\
Multisig health controller; the existing committee & Regional health,
synchronization and delivery observations & Future L1 MSIG, STATUS or
CLEAR; it does not choose a retained height.\\
Cross-region handler; each correct regional replica & Verified public
updates, regional receipts and local task events & Prepare, activate,
request close and transfer responsibilities; it does not replay Force
a second time.\\
Ready and activation gate; each session program & Fixed handoff,
exact anchor, execution state and authority & Waiting or atomic entry
through \cref{alg:bootstrap}; prefetching grants no voting right.\\\bottomrule
\end{tabularx}
\end{table}

Healthy regions keep watching both public streams even while they have
no Active session. They can prepare data and authorities before being
selected. Their public replay cursor, provisional programs and storage
tasks are separate objects. Local lifecycle phases never determine
whether an EXPORT may advance $G$.

\subsection{Soft detection and multisig health control}
\label{sec:runtime-controls}
The Soft Force Detector takes a validated replay checkpoint and the
complete following Espresso prefix. It reproduces the current export
duty and settles its fixed windows as in \cref{sec:windows}. It returns
either no failure, a failed-window descriptor with its first crossing
header and supporting history, or a need for missing evidence. The
descriptor binds the chain, current branch version, export-session pair,
window start, deadline, target and accepted count. Missing namespace or
control evidence is never a zero count.

Only DATA actually appended to $G$ counts. Duplicates, DRAIN, proposal
QCs and invalid EXPORTs do not count. The observed export session can
lag several speculative production sessions. Polling only supplies new
history; its frequency and local reception time do not decide a failure.
Every recipient recomputes a supplied descriptor against its own
authenticated prefix. Protocol~5 remains the sole Force transition:
independent STATUS/CLEAR, the first matching MSIG, otherwise the
automatic check, then the block's EXPORT payload. A computed soft
condition cannot bypass a matching hard intervention in that prelude.

The committee's health policy distinguishes four actions. To exclude
an unsuitable region from future selection, it submits a future OFF.
To end the current target version's production arrangement, it submits
a future MSIG Force, with an OFF if continued exclusion is intended.
After checking synchronization, data service and operational recovery,
it may submit ON for future selections. To remove an automatic cooldown
before expiry, it submits CLEAR explicitly. ON alone does not clear it.
Each operation records its target, authorized changes and prospective
activation time; concrete health thresholds remain an operational
choice under \cref{ass:committee}.

An MSIG targets a branch version. A regular handoff changes $r$ and
therefore does not invalidate it. An earlier Force changes $\beta$ and
can make it stale. Its attached configuration applies only when the
target matches; an independent STATUS follows its own activation rule.
Health observations do not directly change another replica's public
state. All actions pass through finalized L1 and their future Espresso
activation positions under the existing announcement condition.
An ON does not certify future regional synchrony or replace the
regional fault assumptions.

\subsection{Cross-region Handler}
\label{sec:session-manager}
The manager serializes events that modify its local control state.
Data retrieval and state restoration may run concurrently. Its
inputs and outputs have the following meanings.

\begin{table}[htbp]\centering\small
\caption{Session Manager inputs.}
\label{tab:manager-inputs}
\begin{tabularx}{\linewidth}{@{}p{.24\linewidth}Y@{}}\toprule
Event & Validation and routing\\\midrule
\texttt{L1Finalized} & Authenticate the complete control prefix and
authorities; queue future controls without activating them.\\
\texttt{EspressoFinalized} & Authenticate complete ordered history;
consume the state transition produced by \cref{alg:external,alg:force}.\\
\texttt{RegionalFinalized} & Validate the DATA, complete FINAL receipt
and required ancestry; route to the producing pair and offer the result
to the Publish Layer.\\
\texttt{MaterialAvailable} & Check the target pair, anchor's producing
pair, height and hash, requested object and verified result; update only
the matching preparation or fixed handoff.\\
\texttt{CoreEvent} & Accept an event from the named local engine;
record its view, quit, recovery, Completion or retirement progress.\\\bottomrule
\end{tabularx}
\end{table}

One FINAL signature is not a \texttt{RegionalFinalized} event. Existing
participants collect matching signatures into the receipt of
\cref{def:receipt}. A correct holder delivers the complete result to
the manager and the entered source engine, and offers it to the Publish
Layer. The source engine applies the existing receipt and Completion
rules. A program that has not entered only caches the result; receipt
arrival cannot create an engine or skip its handoff checks. Aggregation
is not a new consensus round. $T_{\rm pub}$ still starts when a complete
result becomes available.

The outputs are \texttt{Prepare/Fetch}, naming the target and exact
requested material; \texttt{Activate}, invoking Protocol~4 once;
\texttt{RequestClose}, carrying a verified closure cause;
\texttt{OfferResult}, preserving original DATA and receipt fields;
and \texttt{Persist/Release}, recording state and transferring resource
responsibilities. A cached callback is material, not an activation
command. A late result can be useful to storage without matching the
program's current preparation target.

\begin{protocol}[H]{alg:rotation}{Handle cross-region events and manage sessions}
Each correct regional replica runs this handler independently.
\begin{enumerate}[label=\textbf{\arabic*)},itemsep=3pt]
\item \textbf{Validate the event.} Route authenticated public input,
regional results and local callbacks by their stated context. Preserve
pending requests and source evidence. Reject a purported completion
whose material or session does not match its request.
\item \textbf{Consume public replay.} For new public history, run or
consume \cref{alg:external}, including \cref{alg:force} in its canonical
order. Use its resulting $\beta$, $G$, eligibility snapshots, duty and
boundary records. This handler never increments $\beta$, chooses a
different cutoff or modifies $G$ independently.
\item \textbf{Reconcile handoffs.} From valid terminal receipts,
derive regular boundaries while their version remains open. Consume
fixed Force boundaries from replay. Verify them by
\cref{sec:witness}. Prepare selected local targets; send complete
handoff material to the selected members when available. A future
Force does not block a current regular handoff.
\item \textbf{Reconcile candidates.} Recompute possible future
targets from the current export context, $G$, eligible regions and
known future controls. Update only unfixed preparations. A different
target pair uses a different record. Pause obsolete speculative work;
terminate an unentered program only when verified facts make that
target impossible or assign its fixed handoff to another region.
\item \textbf{Handle closure and results.} For an entered program
whose version is closed, issue the idempotent \texttt{RequestClose}
of \cref{sec:retirement}. Keep its signed routing, historical evidence
and publication duties. Apply terminal receipts to their entered
source engines under the original Completion rule.
\item \textbf{Check and enter.} For each selected local Waiting
program with a fixed handoff, check \cref{def:ready}. In the same local
event, recheck public validity and the durable entry flag, record Ready,
and invoke Protocol~4 Enter atomically. A failed check leaves Waiting
or applies a justified cancellation; it does not partially initialize
the engine.
\item \textbf{Persist and dispatch.} Durably record state changes,
entry and exit records, and required outgoing actions before exposing
them. Dispatch asynchronous work with its full target and anchor
context. Repeated events resume outstanding actions without creating
another logical entry.
\end{enumerate}
\end{protocol}

The public interpreter has one mutation path: Protocol~5 defines
control, window and Force transitions; Protocol~6 authenticates and
applies them and accepts EXPORTs. Protocol~3 consumes those results.
Protocols~1 and~2 update only the named session's regional core.
An implementation may place these modules in one process without
changing this division of authority.

\subsection{Session lifecycle}
\label{sec:session-lifecycle}
The manager maintains the following record for each local program:
\[
\begin{aligned}
\tagg{Program}[\beta,r]=(&\mathit{phase},\mathit{preparing\_plan},
 \mathit{fixed\_handoff},\mathit{anchor\_block},\\
 &\mathit{execution\_state},\mathit{authority\_and\_bounds},
 \mathit{core\_state},\mathit{pending\_fetches},\\
 &\mathit{entered},\mathit{close\_requested},
 \mathit{publication\_and\_retention\_handles}).
\end{aligned}
\]
The six phases describe the local program, not public finality.
\begin{description}[style=nextline,leftmargin=0pt]
\item[Waiting.] The program is preparing a possible or selected
handoff. It may fetch data, recover state and cache evidence. It has
not entered its regional engine, signed a regional vote or started
a regional commit or progress timer.
\item[Ready.] A fixed handoff, its exact anchor state and all entry
dependencies have passed the gate below. Enter follows in the same
serialized event; Ready introduces no extra scheduling delay.
\item[Active.] The program has entered an unfrozen regional
generation. Initial STATUS, the $2\Delta$ wait, NEWVIEW and ANCHOR
certification all occur in Active. The core runs Protocols~1 and~2,
subject only to the explicit retirement interface in
\cref{sec:retirement} when closure is requested.
\item[Frozen.] The current generation has a valid Quit proof.
The original evidence, receipt and recovery rules apply. Force alone
does not establish a Quit proof or cause this phase.
\item[Completed.] The entered engine has obtained a valid terminal
receipt and performed the original Completion action. It has stopped
regional production and cancelled timers as specified there. Data
delivery and retention responsibilities may remain.
\item[Terminated.] The logical program has exited. It processes no
BFT timer and issues no BFT signature. Durable storage and publication
workers retain its required evidence and duties. An exit record prevents
late messages from restarting it.
\end{description}

Normal transitions are Waiting to Ready to Active; Active to Frozen
on Quit; Frozen to Active on the original recovery admission, with
$g$ advanced and $t=0$; and Active or Frozen to Completed on terminal
Completion. Waiting or Ready may terminate when their target is
definitively invalid. Completed may terminate after responsibility
transfer. Active or Frozen may terminate only through the retirement
contract in \cref{sec:retirement}. While close is pending, they retain
their core phase and set \texttt{close\_requested}; there is no seventh
phase.

An unfixed \texttt{preparing\_plan} can change when $G$, the source
duty or candidate eligibility changes. The fixed handoff cannot.
Its target pair, assigned region, initial anchor, source context,
selection position, authority and reference bounds remain immutable.
In particular, a later OFF does not replace a historical regular
selection. Local locks, votes, view entry times and timers belong to
one pair only. Recovering a crashed process resumes this durable
logical record; it does not clear its votes or perform another initial
Enter. Storage and task references carry the full pair and anchor
identity. No additional protocol-level session identifier is needed.

Several regions may hold unfixed Waiting preparations for the same
future pair. Only the region selected by the fixed handoff can enter.
Temporary loss of candidate eligibility may pause work without
terminating a preparation that could become eligible again. A retired
record remains a tombstone even after active resources are released.

An Active old program may already lack permission to extend $G$.
A Completed or Terminated program's existing results may still be
accepted if their version and interval satisfy the public acceptance
rules. The phase does not revoke or create that permission.

\subsection{Handoff records, anchors and Ready}
\label{sec:witness}
Let $\tagg{NextLive}(i,M)$ be the first \emph{other} eligible region
after $i$ in the fixed cyclic order $\mathcal L_R$, searching at most
one cycle. Return $\bot$ if none exists. Eligibility combines the
historical governance flag, cooldown and positive designated-authority
weight. For the synthetic genesis source $\bot$, search from the first
region through the whole list. This version retains the rule requiring
another region.

A pending boundary stores source context
$(i_b,\beta_b,r_b)$ and next index $r_b+1$. A regular boundary uses
its completed predecessor. A Force with an assigned export duty
captures that duty's context. With no assigned duty, a Force reuses
the stored source context. Repeated Forces before assignment replace
the pending cause and version without incrementing an unopened
session. Genesis uses $(\bot,\bot,-1)$.

For REGULAR, the successor of $(\beta,r)$ is $(\beta,r+1)$ and
inherits exactly $T_{\beta,r}$. For Force, the new version's index is
$r_b+1$ and its anchor is exactly $G$ before the triggering block's
payload. The producer of that anchor need not be the captured export
session. The anchor keeps its original producing pair, hash and
commitment evidence.

For a regular boundary set $q_0$ to the height of
$\xi(T_{\beta,r})$; for Force set $q_0$ to its triggering height.
Using the source region just defined and the designated next authority,
the fixed selection is
\[
 q_s=\min\{q\ge q_0:
 \tagg{NextLive}(i_b,\chi_E(q))\ne\bot\},\qquad
 i_{\rm next}=\tagg{NextLive}(i_b,\chi_E(q_s)).
\]
The search uses verified history only. Missing evidence causes a wait,
not a skipped candidate. A later snapshot cannot change an established
minimum. Regular selection need not wait for the terminal DATA to be
published. Its signed selection reference may be historical; no
receive-time freshness rule is added here.

\paragraph{Typed handoff packet.}
A handoff packet contains the fields below. References can name
authenticated cached objects or objects still to be fetched; a reference
alone does not establish Ready.
\begin{table}[htbp]\centering\small
\caption{Handoff packet fields.}
\label{tab:handoff-packet}
\begin{tabularx}{\linewidth}{@{}p{.26\linewidth}Y@{}}\toprule
Field & Type and meaning\\\midrule
Domain and kind & Chain/protocol domain; one of REGULAR, MSIG, AUTO.\\
Source context & Predecessor pair $(\beta_b,r_b)$ and source region
$i_b$, including the stored context when no duty is assigned.\\
Target & Target pair $(\beta',r')$ and selected region $i_{\rm next}$.\\
Anchor & Producing pair, DATA height, hash, full DATA or reference,
and original direct or covering commitment evidence.\\
Authority & Finalized $\tagg{AUTH}_{r',0}$, its fixed configuration,
quota, and publication-height proof.\\
Selection and bounds & $q_0,q_s$, authenticated selection snapshot,
$q^{\rm open}_{\beta',r'}$ and $\lambda^{\rm open}_{\beta',r'}$.\\
Cause and evidence & Mode-specific cause; trusted or validated replay
checkpoint, complete required L1/Espresso history and regional proofs.\\\bottomrule
\end{tabularx}
\end{table}

\paragraph{Handoff verification.}
\label{sec:verify-handoff}
The procedure $\tagg{VerifyHandoff}(C,\mathcal H,P)$ takes an
authorized replay checkpoint $C$, required source history $\mathcal H$
and packet $P$. It returns verified fixed fields, invalid, or needs
evidence. It makes no independent public-state transition.
\begin{enumerate}[label=(\roman*),itemsep=2pt]
\item Check the domain, field types, source and target names. Replay
authenticated compatible history from $C$ through all positions
needed by the cause, selection and currently processed public prefix.
Verify complete applicable controls, not only inclusion of the named
control. Consult validated caches when available.
\item For REGULAR, validate the predecessor's own fixed boundary,
initial authority, exact quota and terminal DATA
$h(T_{\beta_b,r_b})=h(A_{\beta_b,r_b,0,0})+\TB_{r_b}$.
Validate its terminal FINAL receipt, ancestry and signed $\xi$.
Require $\beta'=\beta_b$, $r'=r_b+1$ and that version still open in
the processed prefix. The anchor must equal the terminal DATA.
\item For MSIG, verify its authorized finalized L1 record, control
sequence and first-crossing header $J_{q(\theta_a)}$. Replay the
entire prelude and check that this is the first due Force whose target
matches the entry version. Check attached changes only under this
match. A stale instruction or another due instruction consumed by the
same prelude cannot form a second boundary.
\item For AUTO, replay the fixed window's count and first deadline
crossing from complete history. Check insufficient accepted DATA and
the absence of a matching higher-priority MSIG in that prelude.
Check the cooldown derived at the realized triggering timestamp.
A supplied alarm or a local timeout is not proof.
\item For either Force, obtain the captured source context, old and
new versions, retained prefix and boundary from that unique replay
transition. Require $\beta'=\beta_{\rm entry}+1$, $r'=r_b+1$ and
anchor equality with pre-payload $G$, including prefix provenance.
During an unassigned interval, verify the stored context and any
intervening Forces rather than inventing an empty export session.
Require the resulting target version still open in the processed
prefix. An intervening effective Force supersedes this boundary.
\item Recompute the first eligible $q_s$ from $q_0$ and all required
historical snapshots. Check target region and immutable next authority.
Set $q^{\rm open}_{\beta',r'}=q_s$ and
$\lambda^{\rm open}_{\beta',r'}$ to the maximum of the anchor's L1
reference height and the initial authority's publication height.
Validate the anchor's original direct or covering commitment evidence
under \cref{sec:coverage}, its complete block and required ancestry.
Do not relabel it with the target pair.
\item Compare the derived fixed fields with $P$ and any existing
fixed record for its target pair. Equal duplicates share that record;
different fixed fields for the same pair are rejected. A packet cannot
replace an already installed handoff. Before entry, recheck validity
against the manager's current processed public prefix.
\end{enumerate}

A Force without an eligible successor still has a fixed cause,
anchor, source context and new version. The interpreter retains that
pending boundary. It has no complete selected packet and cannot make
a program Ready until the first eligible historical position fixes
the target region. A later Force can supersede it while preserving the
same stored predecessor context. The synthetic genesis boundary is
verified against the initial checkpoint and selection rule of
\cref{sec:l1}; it does not claim a nonexistent REGULAR, MSIG or AUTO
cause.

\begin{definition}[\defterm{Ready}]
\label{def:ready}
A local Waiting program for $(\beta,r)$ is Ready precisely when:
\begin{enumerate}[label=(\roman*)]
\item handoff verification fixes this target pair and selects the
replica's region, or the authenticated genesis rule does so;
\item its exact anchor's original producing pair, height, hash,
complete DATA and commitment evidence have been validated;
\item the native execution state needed to continue after this
anchor has been recovered and verified against that same anchor;
\item the matching finalized authority, fixed quota, separate
reference lower bounds and all validation dependencies are available;
\item the manager's processed public history has not closed or
superseded the handoff; and
\item the durable record has never performed initial Enter and has
no Terminated tombstone.
\end{enumerate}
\end{definition}

State recovery uses the application's existing validation interface;
it does not add a new execution-trust assumption. A cached state for
DATA 110 cannot make a program with fixed anchor 112 Ready. A future
Force announcement and fully downloaded candidate state still leave
the program Waiting. Each selected correct member relays the complete
validated bundle when it passes the gate. Final gate checks and Enter
are one serialized local event with respect to public replay.

\subsection{Activation and bootstrap}
\label{sec:ready}
\begin{protocol}[H]{alg:bootstrap}{Enter and bootstrap the selected session}
Replica $v\in V_{i_{\beta,r}}$ runs these steps for one fixed target
$(\beta,r)$. Let $\Delta=\Delta_{i_{\beta,r}}$.
\begin{enumerate}[label=\textbf{\arabic*)},itemsep=3pt]
\item \textbf{Enter.} In the Ready event, atomically check and set
the durable initial-entry flag. Install $A_{\beta,r,0,0}$,
$\tagg{AUTH}_{r,0}$, $\TB_r$ and both opening reference bounds.
Set the entry lock to $Q^0_{\beta,r}$. Enter $(\beta,r,0,0)$
unfrozen, change Ready to Active, record $e_v(\beta,r,0,0)$ once,
and start the original progress checks. Send STATUS with this lock
and the fixed handoff to the initial leader, and process buffered
proofs for this session. Persist entry before dispatching signatures.
\item \textbf{Choose the anchor.} The initial leader waits
$2\Delta$ from its own entry. If the view remains current and
unfrozen and the retirement contract permits the action, send NEWVIEW
choosing exactly $A_{\beta,r,0,0}$, justified by $Q^0_{\beta,r}$
and the fixed handoff.
\item \textbf{Certify and produce.} Members apply the original
ANCHOR-vote step of \cref{alg:recovery}. On the ANCHOR QC, the leader
starts \cref{alg:steady} at DATA height $h(A_{\beta,r,0,0})+1$
and physical position 1. Local failures use the same recovery rules.
All these steps occur in Active, except when the original Quit,
Completion or explicit retirement transition intervenes.
\end{enumerate}
\end{protocol}

Let $e=\min_v e_v(\beta,r,0,0)$ over correct members. If the fixed
handoff remains valid and no retirement interrupts bootstrap, the
complete-bundle relay and existing regional communication bound make
all correct members enter by $e+\Delta$. With a correct initializer
and no quit, NEWVIEW is sent by $e+3\Delta$, its ANCHOR QC reaches
the leader by $e+5\Delta$, and the first DATA proposal reaches all
correct members by $e+6\Delta$. Ready and Enter are in the same local
event, so these are the original entry-based bounds. No asynchronous
Ready queue or additional scheduling allowance is introduced.
State transfer before entry uses direct transfer or authenticated
Espresso retrieval under the existing assumptions. The bounds do not
start at the beginning of transfer.

\paragraph{The public export clock.}
Public replay assigns an export duty independently of a member's
phase. Its window starts at the containing EXPORT's timestamp for
an already selectable normal successor, at the triggering block's
timestamp for an immediately selected forced successor, or at the
first eligible selection block when the boundary previously had no
assigned duty. A selected program may still be Waiting while that
public window runs. Download completion, Ready, process restart and
duplicate handoffs never reset its deadline. Only an unassigned duty
has no DATA window. Missing evidence pauses replay at its canonical
position; it does not create a later public opening time.

\subsection{Concurrent candidates and races}
\label{sec:runtime-races}
Suppose A's session is $(7,20)$ with terminal DATA 110, B already
produces in $(7,21)$, and $G$ ends at 108. A future MSIG targets
version 7. B can retain its old Active program while preparing
candidate $(8,21)$ from 108. If $G$ advances to 109, that preparation
can change to 109; the old Active program's initial anchor remains
110. If Force then fixes B as successor with anchor 109, all further
preparation is for that fixed anchor. B can enter $(8,21,0,0)$ while
$(7,21)$ closes under \cref{sec:retirement}. Their locks, votes,
entry times and timers are distinct.

If A's terminal 110 enters $G$ before Force, the export context can
advance to B. The candidate then becomes $(8,22)$ at the next eligible
region after B. With cyclic order A, B, C, D, a prospective OFF(C)
that has become effective before the Force selection can redirect
this still-unfixed candidate to D. If Force finally retains B's 112,
the fixed target is D's $(8,22)$ with original anchor producer
$(7,21)$. D must verify state at 112 before starting DATA 113.
Earlier preparation of $(8,21)$ or of state at 109 grants no entry.
OFF cannot rewrite an already fixed regular handoff.

A terminal receipt received before Force can prepare or start a
regular successor in the old version. Once the same node verifies
Force, it cancels an invalid unentered target or requests retirement
of the entered one. If Force arrives first, the late receipt is
retained as regional evidence and cannot resurrect the closed target.
Nodes can briefly have different speculative programs; the same
verified public prefix fixes the same public state and Force boundary.

Every fetch callback retains its original target and anchor identity.
A reply for $(7,21)$ at 110 cannot make $(8,21)$ at 108 Ready.
An old candidate reply may be cached and independently validated for
another use, but it cannot overwrite a fixed anchor. Duplicate receipts,
packets and callbacks cannot run initial Enter twice. A stale MSIG
following an earlier AUTO neither closes the new version nor applies
its attached configuration. Reconciliation discards its obsolete
preparation basis without invalidating a different fixed handoff.

\subsection{Graceful close and termination}
\label{sec:runtime-close}
Public closure and program exit are separate events. Protocol~5
immediately removes a closed version's permission for new global
appends. Protocol~3 then delivers \texttt{RequestClose} to its local
entered programs, using the verified closure cause. The request is
idempotent and records \texttt{close\_requested}; it is neither Quit
nor terminal Completion. The exact restrictions on new work, treatment
of existing timers and messages, and admissible \texttt{FinishClose}
condition are defined by \cref{sec:retirement}. This manager interface
adds no shutdown-time assumption or fixed retirement deadline. A new
session need not wait for the old program to finish closing.

For an unentered program, definitive invalidation cancels its pending
entry and tasks and records Terminated. It has no regional timers or
votes to drain. For Completed, the original Completion action has
already stopped the core. For an Active or Frozen program, only the
retirement contract can authorize \texttt{FinishClose}. An empty local
timer list, a different current public version, or a locally observed
terminal block alone is not that authorization.

In every case resource release requires durable exit and responsibility
records. Required committed DATA, receipts, regional evidence, anchors
and recoverable state must remain accessible under the existing
retention rules. Outstanding complete results and publication retries
are transferred to the Publish Layer, and pending requests in discarded
suffixes remain available for reproposal. Storage and publication
workers acknowledge the responsibilities they accept before the
program releases its last required copy. Later messages can update
those services without reactivating the logical engine. Terminated
does not undo regional commitments, remove DATA from $G$, or authorize
stake withdrawal.
